\section{引言}
\label{sec:intro}

技术报告的价值在于可核对：读者能顺着数字找到实验，顺着实验找到代码。排版不是
这份价值的来源，但它决定读者愿不愿意读下去，也决定作者每写一份报告要不要把
同样的决定重做一遍。

我们具体被三件事绊住过。第一，每次新开报告都要从空白文档重调页边距、字号、
图注和表格线，而调出来的东西每份都不太一样。第二，中英混排反复踩坑：中文要
XeLaTeX 才能正确断行，英文用 pdfLaTeX 最省事，而两套环境下的观感必须一致。
第三，数字写法没有统一约定，同一张表里出现 \texttt{62.4} 和 \texttt{62.40}，
读者就有理由怀疑哪一个是笔误。

这份模板针对的就是这三件事，贡献可以归纳为三点：

\begin{contributions}
  \item 一个类文件 \texttt{boogonreport.cls} 承载全部版式决策：标题块、摘要、
        题注、三线表、代码块、参考文献与附录；
  \item 一个 \texttt{zh} 选项同时切换字体管线与本地化文案，中文报告不再需要
        手写一串 \verb|\renewcommand|；
  \item 一套构建脚本与 CI，让任何一台装了 TeX Live 的机器上，一条命令就能得到
        同一份 PDF。
\end{contributions}

全文组织如下：\secref{sec:layout} 说明版面参数及其来源；\secref{sec:type}
交代中英两条字体管线；\secref{sec:figure} 给出图表与数字的写法；
\secref{sec:build} 是构建与发布流程；\secref{sec:limit} 列出这套模板不打算
解决的问题。
